% GENERATED FILE. Edit canonical lessons, then run npm run generate:books.
% canonical-source-hash: fnv1a64:b12cae49fe7715e5
% canonical-lessons: RU-C152-okolo, RU-C152-vokrug, RU-C152-vnutri, RU-C152-snaruzhi, RU-C152-szadi

\chapter{Beside, Around, Inside, Outside, Behind}
\label{ch:ru-beside-around-inside-outside-behind}
\hlchaptermodality{152}

\begin{chapteropening}
\textbf{By the end of this chapter:} \emph{I can say beside, around, inside, outside and behind.}

Complete the last lesson of chapter 152: I can say beside, around, inside, outside and behind.
\end{chapteropening}

\section[\texorpdfstring{\ru{около} (ókolo)}{около (ókolo)}]{\ru{около} (ókolo) — beside; roughly (with numbers)}

\label{lesson:RU-C152-okolo}



\begin{quote}
Before the new one: say the Russian for a district, then the Russian for an avenue.
\end{quote}

\subsection*{You'll want to know: \ru{около}}
\textbf{\ru{около}} (\emph{ókolo}) — \textquotedblleft{}beside; roughly (with numbers)\textquotedblright{}.

\begin{cousinweb}
\textbf{\ru{около}} is \emph{\ru{о}-}, \textquotedblleft{}around\textquotedblright{}, with \emph{\ru{коло}}, \textquotedblleft{}a circle, a wheel\textquotedblright{} (Vasmer), the root of \emph{\ru{колесо}}, \textquotedblleft{}wheel\textquotedblright{}: \textquotedblleft{}round about\textquotedblright{}, and so \textquotedblleft{}approximately\textquotedblright{}.
\end{cousinweb}

\begin{grammarlens}[title={What you've built}]
One more word for this chapter.
\end{grammarlens}

\subsection*{Your turn}
\begin{itemize}
\raggedright
  \item \emph{Say it:} \emph{ókolo}
  \item \emph{Say it:} \emph{ókolo}, once more
  \item \emph{Say it:} say \emph{ókolo}
  \item \emph{Recall:} say the Russian for a district, then the Russian for an avenue, then say \emph{ókolo} again
\end{itemize}

\begin{tcolorbox}[breakable,colback=teal!4,colframe=teal!35!black,title={Before you move on}]
What does \ru{около} mean? (\textquotedblleft{}Beside; roughly (with numbers)\textquotedblright{}.) Say it once more.
\end{tcolorbox}
\section[\texorpdfstring{\ru{вокруг} (vokrúg)}{вокруг (vokrúg)}]{\ru{вокруг} (vokrúg) — around}

\label{lesson:RU-C152-vokrug}



\begin{quote}
Before the new one: say the Russian for an avenue, then the Russian for beside; roughly.
\end{quote}

\subsection*{You'll want to know: \ru{вокруг}}
\textbf{\ru{вокруг}} (\emph{vokrúg}) — \textquotedblleft{}around\textquotedblright{}.

\begin{cousinweb}
\textbf{\ru{вокруг}} is \emph{\ru{в}}, \textquotedblleft{}in\textquotedblright{}, with \emph{\ru{круг}}, \textquotedblleft{}a circle\textquotedblright{}: \textquotedblleft{}in a ring around\textquotedblright{}.
\end{cousinweb}

\begin{grammarlens}[title={What you've built}]
One more word for this chapter.
\end{grammarlens}

\subsection*{Your turn}
\begin{itemize}
\raggedright
  \item \emph{Say it:} \emph{vokrúg}
  \item \emph{Say it:} \emph{vokrúg}, once more
  \item \emph{Say it:} say \emph{vokrúg}
  \item \emph{Recall:} say the Russian for an avenue, then the Russian for beside; roughly, then say \emph{vokrúg} again
\end{itemize}

\begin{tcolorbox}[breakable,colback=teal!4,colframe=teal!35!black,title={Before you move on}]
What does \ru{вокруг} mean? (\textquotedblleft{}Around\textquotedblright{}.) Say it once more.
\end{tcolorbox}
\section[\texorpdfstring{\ru{внутри} (vnutrí)}{внутри (vnutrí)}]{\ru{внутри} (vnutrí) — inside}

\label{lesson:RU-C152-vnutri}



\begin{quote}
Before the new one: say the Russian for beside; roughly, then the Russian for around.
\end{quote}

\subsection*{You'll want to know: \ru{внутри}}
\textbf{\ru{внутри}} (\emph{vnutrí}) — \textquotedblleft{}inside\textquotedblright{}.

\begin{cousinweb}
\textbf{\ru{внутри}} is \emph{\ru{в}}, \textquotedblleft{}in\textquotedblright{}, with a form of \emph{\ru{нутро}}, \textquotedblleft{}the inside\textquotedblright{}, from \emph{*vŭn-ǫtrĭ}, \textquotedblleft{}within\textquotedblright{} (Vasmer). Vasmer relates it to Greek \emph{éntera}, \textquotedblleft{}entrails\textquotedblright{}, and Latin \emph{interus}, \textquotedblleft{}inner\textquotedblright{}.
\end{cousinweb}

\begin{grammarlens}[title={What you've built}]
One more word for this chapter.
\end{grammarlens}

\subsection*{Your turn}
\begin{itemize}
\raggedright
  \item \emph{Say it:} \emph{vnutrí}
  \item \emph{Say it:} \emph{vnutrí}, once more
  \item \emph{Say it:} say \emph{vnutrí}
  \item \emph{Recall:} say the Russian for beside; roughly, then the Russian for around, then say \emph{vnutrí} again
\end{itemize}

\begin{tcolorbox}[breakable,colback=teal!4,colframe=teal!35!black,title={Before you move on}]
What does \ru{внутри} mean? (\textquotedblleft{}Inside\textquotedblright{}.) Say it once more.
\end{tcolorbox}
\section[\texorpdfstring{\ru{снаружи} (snarúzhi)}{снаружи (snarúzhi)}]{\ru{снаружи} (snarúzhi) — outside}

\label{lesson:RU-C152-snaruzhi}



\begin{quote}
Before the new one: say the Russian for around, then the Russian for inside.
\end{quote}

\subsection*{You'll want to know: \ru{снаружи}}
\textbf{\ru{снаружи}} (\emph{snarúzhi}) — \textquotedblleft{}outside\textquotedblright{}.

\begin{cousinweb}
\textbf{\ru{снаружи}} is \emph{\ru{с}}, \textquotedblleft{}from\textquotedblright{}, with \emph{\ru{наружи}}, \textquotedblleft{}outside\textquotedblright{}, built on dialect \emph{\ru{ружь}}, \textquotedblleft{}the outward look\textquotedblright{} (Vasmer), which he relates to Latvian \emph{raũdzît}, \textquotedblleft{}to look\textquotedblright{}.
\end{cousinweb}

\begin{grammarlens}[title={What you've built}]
One more word for this chapter.
\end{grammarlens}

\subsection*{Your turn}
\begin{itemize}
\raggedright
  \item \emph{Say it:} \emph{snarúzhi}
  \item \emph{Say it:} \emph{snarúzhi}, once more
  \item \emph{Say it:} say \emph{snarúzhi}
  \item \emph{Recall:} say the Russian for around, then the Russian for inside, then say \emph{snarúzhi} again
\end{itemize}

\begin{tcolorbox}[breakable,colback=teal!4,colframe=teal!35!black,title={Before you move on}]
What does \ru{снаружи} mean? (\textquotedblleft{}Outside\textquotedblright{}.) Say it once more.
\end{tcolorbox}
\section[\texorpdfstring{\ru{сзади} (szádi)}{сзади (szádi)}]{\ru{сзади} (szádi) — behind}

\label{lesson:RU-C152-szadi}



\begin{quote}
Before the new one: say the Russian for inside, then the Russian for outside.
\end{quote}

\subsection*{You'll want to know: \ru{сзади}}
\textbf{\ru{сзади}} (\emph{szádi}) — \textquotedblleft{}behind\textquotedblright{}.

\begin{cousinweb}
\textbf{\ru{сзади}} is \emph{\ru{с}}, \textquotedblleft{}from\textquotedblright{}, with \emph{\ru{зади}}, from \emph{\ru{зад}}, \textquotedblleft{}the back, the rear\textquotedblright{}.
\end{cousinweb}

\begin{grammarlens}[title={What you've built}]
One more word for this chapter.
\end{grammarlens}

\subsection*{Your turn}
\begin{itemize}
\raggedright
  \item \emph{Say it:} \emph{szádi}
  \item \emph{Say it:} \emph{szádi}, once more
  \item \emph{Say it:} say \emph{szádi}
  \item \emph{Recall:} say the Russian for inside, then the Russian for outside, then say \emph{szádi} again
\end{itemize}

\begin{tcolorbox}[breakable,colback=teal!4,colframe=teal!35!black,title={Before you move on}]
What does \ru{сзади} mean? (\textquotedblleft{}Behind\textquotedblright{}.) Say it once more.
\end{tcolorbox}
